\documentclass[11pt]{article}

\usepackage[margin=25mm]{geometry}

\usepackage{amsmath,amssymb,fontspec,unicode-math}

\IfFontExistsTF{DejaVu Serif}{\setmainfont{DejaVu Serif}}{\setmainfont{Latin Modern Roman}}

\IfFontExistsTF{Latin Modern Math}{\setmathfont{Latin Modern Math}}{\setmathfont{DejaVu Math TeX Gyre}}

\usepackage{xurl,hyperref}

\setlength{\parindent}{0pt}\setlength{\parskip}{6pt}

\begin{document}

\title{Synthetic Series-Symbol Data Generation for Time Series Foundation Models}

\author{ELVIA preliminary AI analysis}\date{21 October 2026}\maketitle

Preliminary assessment: Minor.\par

Independent scientific review remains pending. Author adoption has not been requested.

Printed stationarity and corpus-diagnostic inferences need local correction; current release paths also expose test-based checkpoint selection and a padding loss-mask discrepancy.

\section*{Review scope}

Full supplied technical report read alongside the corresponding original-paper claim passages and applicable author-release source. This pass independently checks the stated reasoning and bounds; historical executions were not rerun.

\section*{Problem 1: Printed AR conditions do not ensure causal stationarity}

Source: Section 3.1.2, Eq.(1), Sampling Strategy, PDF p.4

Assessment: supported; impact: minor

\subsection*{Source claim (paraphrase)}

The paper says coefficients in \((-1,1)\), \(\sum_i\phi_i<1\) and \(|\phi_p|<1\) enforce stationarity of its ARMA generator.

\subsection*{Analysis}

Take \(\phi_1=-9/10,\phi_2=9/10,\theta_1=1/5\). Both printed tests hold, but the companion polynomial factors as \((\lambda-3/5)(\lambda+3/2)\), with unstable eigenvalue \(-3/2\). For the stated MA sign, the causal innovation response is \(\psi_k=(4/21)(3/5)^k+(17/21)(-3/2)^k\), whose squared sum diverges. The MA factor does not cancel the unstable root.

\subsection*{Finding and scope}

The manuscript’s sufficient conditions are false for a causal finite-variance recursion. Current S2Generator’s default stationary mode constructs stable roots instead, so this does not show that its current generated inputs are unstable or identify the published corpus’s generator revision.

\subsection*{Proposed correction}

Require all uncancelled AR lag roots outside the closed unit disk, or conservatively require \(\sum_i|\phi_i|<1\). Exact stationarity also needs the appropriate initial law.

The described correction remains a proposal; author adoption or a corresponding erratum is not established.

\section*{Problem 2: Mean ADF p-value does not identify nonstationary prevalence}

Source: Appendix B.3, Table 10, current arXiv v3 physical PDF p.44, printed p.26

Assessment: supported; impact: minor

\subsection*{Source claim (paraphrase)}

The paper interprets an average ADF p-value above \(0.05\), including \(0.0538\), as suggesting most samples are nonstationary.

\subsection*{Analysis}

A set of \(1000\) p-values with \(940\) equal to \(0.001\) and \(60\) equal to \(0.881\) has mean exactly \(0.0538\), although \(94\%\) are below the rejection threshold. Moreover, failure to reject a unit-root null is not proof of nonstationarity. The table may correctly report its averages; those averages do not determine either the rejection fraction or the underlying prevalence.

\subsection*{Finding and scope}

The aggregate-to-majority inference is unsupported. The actual per-series results and corpus nonstationarity fraction remain unknown.

\subsection*{Proposed correction}

Report individual test specifications and rejection indicators, then summarize the nonrejection fraction with sampling uncertainty, without interpreting nonrejection as proof of the null.

The described correction remains a proposal; author adoption or a corresponding erratum is not established.

\section*{Problem 3: Spectral concentration is not a certificate of nonchaotic dynamics}

Source: Appendix B.3 and B.5, forecastability interpretation

Assessment: supported; impact: minor

\subsection*{Source claim (paraphrase)}

Forecastability above \(0.3\) is used to conclude that generated outputs are not chaotic and are predictable.

\subsection*{Analysis}

Consider \(F(x,\theta)=(2x\bmod1,\theta+\pi/4\bmod2\pi)\) and observation \(y_t=\cos\theta_t+\varepsilon(x_t-1/2)\). The expanding coordinate retains Lyapunov exponent \(\log2>0\). For a length-\(256\) sample, the cosine occupies one Fourier bin; as positive \(\varepsilon\) tends to zero, normalized spectral-entropy forecastability tends to \(1\). A dominant periodic observed component therefore does not remove chaos in the generating system.

\subsection*{Finding and scope}

The claimed universal diagnostic implication fails. This example does not establish that the actual synthetic corpus is chaotic, nor deny usefulness of the spectral feature as a descriptive statistic.

\subsection*{Proposed correction}

Limit the conclusion to spectral concentration under the declared estimator. Claims about chaos or forecast accuracy need separately justified dynamics or out-of-sample forecasting tests.

The described correction remains a proposal; author adoption or a corresponding erratum is not established.

\section*{Problem 4: Classification and M4 checkpoint selection uses test targets}

Source: Table 15; SymTime exp/exp\_classification.py, train; exp/exp\_short\_term\_forecasting.py; data\_provider/data\_loader.py, Dataset\_M4

Assessment: supported; impact: minor

\subsection*{Source claim (paraphrase)}

The released classification and short-term forecasting paths use the official test targets for their checkpoint-selection criterion.

\subsection*{Analysis}

Classification training requests TRAIN, TEST and TEST; validation accuracy on the first TEST call drives early stopping, and the final evaluation uses that same split. For M4, the validation path reads the nontraining archive and passes those future targets into the loss selecting the checkpoint. Thus the selected predictor depends on test labels through model selection, even without their entering gradient updates.

\subsection*{Finding and scope}

The current release does not yield an independently selected fresh-test estimate on those paths. Published run-to-code provenance, effect sizes and corrected rankings are not established.

\subsection*{Proposed correction}

Split validation data from training histories or training examples, select all epochs/hyperparameters there, and reserve official test targets for evaluation after selection.

The described correction remains a proposal; author adoption or a corresponding erratum is not established.

\section*{Problem 5: Masked-time loss includes already invalid padding}

Source: Eq.(2), PDF p.5; models/pretrain\_model.py, lines 158–180 and 250–270

Assessment: supported; impact: minor

\subsection*{Source claim (paraphrase)}

The stated loss reconstructs randomly masked real patches, but the release averages every position false in the postmask attention mask.

\subsection*{Analysis}

Let original validity be \(v=(1,1,0,0)\), postmask validity \(a=(0,1,0,0)\), true patches \((2,3,0,0)\), and predictions \((0,3,1,1)\). Squared errors are \((4,0,1,1)\). Using \(1-a\) gives loss \(2\), whereas the intended real-mask selector \(v(1-a)\) gives \(4\). Padding gradients are spuriously nonzero. Even zero padded residuals still dilute the denominator. If every input patch is valid, the two selectors coincide.

\subsection*{Finding and scope}

The loss-mask discrepancy is concrete for padded inputs. The original serialized pretraining masks and frequency of affected examples are unknown, so no quantitative downstream damage is inferred.

\subsection*{Proposed correction}

Use the original-validity intersection \(v(1-a)\) for both numerator and denominator, with an explicit policy for no masked valid patches.

The described correction remains a proposal; author adoption or a corresponding erratum is not established.

\section*{Limits}

Preliminary AI source/argument assessment; independent human scientific adjudication remains pending.\par

No author-provided or retained verification program was executed. Published empirical results were not independently reproduced.\par

\section*{Source identity}

Authors’ arXiv:2510.08445v3, current official version checked 2026-10-21

\url{https://arxiv.org/pdf/2510.08445v3}

Reviewed PDF SHA-256: \nolinkurl{b7d1b9da8a7613a76b3dc76a10f337ec2109c1a42e8ce9e315d99f7ec11a0813}

The SHA-256 identifies the exact current-author-version PDF inspected. The original acquisition timestamp for this retained PDF is unknown; file modification times were not used. Current version identity was independently checked against official arXiv metadata.

\end{document}
